% Provenance (never printed): see sources/outline_sources.md (Year 2: tholeiitic, calc-alkaline
% and alkaline series; the TAS diagram; settings; isotopic signatures of mantle and crust).
% Classification from BGS RR 99-06 (after Streckeisen 1976, Le Bas et al. 1986, Le Maitre et
% al. 1989); series and data from Vermeesch and Pease (2021) and the GeoplotR package;
% curves and counts from figdata/bachelor-2/series.py; exercise and problem values in the
% Book 4 scratch calculators calc/ch09.py and calc/ch09pb.py.
\chapter{Magmatic Series and Tectonic Settings}\label{ch:b2:magmatic-series}

Hand a petrologist a chemical analysis and no map, and she can often say whether the
lava erupted at a ridge, on an ocean island or above a subducting plate. The lavas
of Iceland and those of the Cascade volcanoes of western North America both run
from basalt to rhyolite; analysed by the hundred, they follow two different paths as
they evolve, one growing ever richer in iron, the other not. This chapter first
names igneous rocks, from their minerals when they can be seen and from their
chemistry when they cannot; it then groups them into series, families of rocks
related by differentiation, and ties each series to a tectonic setting, with the
diagrams that petrologists use to read a setting from an analysis, and their limits.
% ledger: vp21:series

\begin{recall}
\textbf{From the high-school volume}: basalt, andesite, rhyolite, gabbro and granite;
felsic and mafic rocks; volcanic and plutonic textures. \textbf{From the Year 1
volume}: mid-ocean ridges, subduction zones and volcanic arcs, hotspots and large
igneous provinces, continental rifts. \textbf{From \cref{ch:b2:elements-in-the-earth}}:
MORB and OIB, primitive-mantle-normalised (spider) diagrams, compatible and
\omterm{def:b2:elements-in-the-earth:partition}{incompatible elements}. \textbf{From \cref{ch:b2:isotope-geochemistry}}: the \omterm{def:b2:isotope-geochemistry:reservoirs}{depleted
mantle} and the enriched isotopic reservoirs. \textbf{From
\cref{ch:b2:magmatic-differentiation}}: \omterm{def:b2:magmatic-differentiation:crystallisation}{fractional crystallisation}, the \omterm{def:b2:magmatic-differentiation:lld}{liquid line of
descent}, \omterm{def:b2:magmatic-differentiation:harker}{Harker diagrams} and the Rayleigh law.
\end{recall}

\section{Naming igneous rocks}\label{sec:b2:magmatic-series:naming}

A coarse-grained rock shows its minerals, and is named from them. A lava is often too
fine-grained or too glassy for that, and is named from its chemical analysis. Both
schemes are international conventions: they describe, and imply nothing about how
the rock formed.
% ledger: bgs:tas-use

\begin{definition}[Modal composition; QAPF diagram]\label{def:b2:magmatic-series:qapf}
The \emph{modal composition}\index{modal composition} of a rock is the proportion by
volume of each of its minerals. The \emph{QAPF diagram}\index{QAPF diagram} names a
coarse-grained igneous rock with less than \qty{90}{\percent} of mafic minerals from
four modal components recalculated to a sum of 100: quartz (Q), alkali feldspar
(A), plagioclase (P) and the feldspathoids (F). Since quartz and feldspathoids do not
crystallise together, every rock plots in one of two triangles, QAP or FAP, joined
along their common side AP.
\end{definition}
% ledger: bgs:qapf

\begin{omfigure}
\centering
\begin{tikzpicture}[font=\scriptsize, x=1cm, y=1cm]
\def\L{6.4}
\pgfmathsetmacro{\H}{\L*0.8660254}
% upper point (q, p): q = Q fraction, p = P/(A+P)
\newcommand{\qapfup}[2]{({(#1)*\L/2 + (#2)*(1-(#1))*\L}, {(#1)*\H})}
\newcommand{\dn}[2]{({(#1)*\L/2 + (#2)*(1-(#1))*\L}, {-(#1)*\H})}
\draw[thick] (0,0) -- (\L/2,\H) -- (\L,0) -- (\L/2,-\H) -- cycle;
\draw (0,0) -- (\L,0);
\foreach \q in {0.9,0.6,0.2,0.05} \draw \qapfup{\q}{0} -- \qapfup{\q}{1};
\foreach \q in {0.1,0.6} \draw \dn{\q}{0} -- \dn{\q}{1};
\foreach \p in {0.1,0.65,0.9} \draw \qapfup{0.6}{\p} -- \qapfup{0}{\p};
\draw[dashed] \qapfup{0.6}{0.35} -- \qapfup{0.2}{0.35};
\draw \qapfup{0.2}{0.35} -- \qapfup{0}{0.35};
\foreach \p in {0.1,0.35,0.65,0.9} \draw \dn{0}{\p} -- \dn{0.1}{\p};
\foreach \p in {0.1,0.5,0.9} \draw \dn{0.1}{\p} -- \dn{0.6}{\p};
\node[anchor=south] at (\L/2,\H) {\normalsize Q};
\node[anchor=east] at (0,0) {\normalsize A};
\node[anchor=west] at (\L,0) {\normalsize P};
\node[anchor=north] at (\L/2,-\H) {\normalsize F};
\node at \qapfup{0.94}{0.5} {1a};
\node at \qapfup{0.75}{0.5} {1b};
\node[font=\tiny] at \qapfup{0.4}{0.05} {2};
\node at \qapfup{0.4}{0.225} {3a};
\node at \qapfup{0.4}{0.5} {3b};
\node at \qapfup{0.4}{0.775} {4};
\node[font=\tiny] at \qapfup{0.4}{0.95} {5};
\foreach \p/\n in {0.225/7*,0.5/8*,0.775/9*} \node at \qapfup{0.125}{\p} {\n};
\node[font=\tiny] at \qapfup{0.125}{0.05} {6*};
\node[font=\tiny] at \qapfup{0.125}{0.95} {10*};
\foreach \p/\n in {0.225/7,0.5/8,0.775/9} \node at \qapfup{0.025}{\p} {\n};
\foreach \p/\n in {0.225/7',0.5/8',0.775/9'} \node at \dn{0.05}{\p} {\n};
\node[font=\tiny] at \dn{0.3}{0.05} {11};
\node at \dn{0.35}{0.3} {12};
\node at \dn{0.35}{0.7} {13};
\node[font=\tiny] at \dn{0.3}{0.95} {14};
\node at \dn{0.75}{0.5} {15};
\node[anchor=east] at \qapfup{0.9}{0} {90};
\node[anchor=east] at \qapfup{0.6}{0} {60};
\node[anchor=east] at \qapfup{0.2}{0} {20};
\node[anchor=east] at \dn{0.1}{0} {10};
\node[anchor=east] at \dn{0.6}{0} {60};
\foreach \p/\n in {0.1/10,0.35/35,0.65/65,0.9/90} \node[anchor=south, fill=white, inner sep=0.5pt] at \qapfup{0.6}{\p} {\tiny\n};
\node[anchor=west, align=left, font=\scriptsize] at (\L+0.9,1.6) {Upper triangle: \% Q\\in Q + A + P;\\
  lines of constant\\$100\,P/(A+P)$\\at 10, 35, 65, 90};
\node[anchor=west, align=left, font=\scriptsize] at (\L+0.9,-1.7) {Lower triangle: \% F\\in F + A + P;\\
  $100\,P/(A+P)$\\at 10, 50, 90};
\end{tikzpicture}
\par\smallskip{\small The QAPF classification of coarse-grained igneous rocks. 1a
quartzolite, 1b quartz-rich granitoid, 2 alkali-feldspar granite, 3a syenogranite and
3b monzogranite (together granite), 4 granodiorite, 5 tonalite; 6*--10* the
quartz-bearing varieties of 6--10: 6 alkali-feldspar syenite, 7 syenite, 8 monzonite, 9 monzodiorite
or monzogabbro, 10 diorite, gabbro or anorthosite; 6'--10' their foid-bearing
varieties; 11 foid syenite, 12 foid monzosyenite, 13 foid monzodiorite or
monzogabbro, 14 foid diorite or gabbro, 15 foidolite. Thin fields 6, 6', 10 and 10'
along the margins are not numbered. After the British Geological Survey scheme (based
on Streckeisen 1976).}
\end{omfigure}
% ledger: bgs:qapf

Field 10 holds three names. The \omterm{def:b2:magmatic-series:qapf}{QAPF diagram} does not separate diorite, gabbro and
anorthosite, nor the gabbroic rocks among themselves; further diagrams, based on the
plagioclase and the mafic minerals, do that.

\begin{definition}[Total alkali--silica diagram]\label{def:b2:magmatic-series:tas}
The \emph{total alkali--silica diagram}\index{total alkali--silica diagram} (TAS
diagram) names a volcanic rock whose mode cannot be measured from its mass
percentages of \ce{Na2O + K2O} against \ce{SiO2}, after the analysis has been
recalculated to \qty{100}{\percent} free of water and carbon dioxide. It is divided
into fifteen fields by straight lines whose intersections have fixed coordinates
(the figure after the next paragraph).
\end{definition}
% ledger: bgs:tas, bgs:tas-use

The diagram has three groups of fields. Along the bottom, at low alkalis, lie the
names of the common series, picrobasalt and basalt to rhyolite, separated at 45, 52,
57 and \qty{63}{\percent} silica and, for rhyolite, by a sloping line. Above them, the S fields and trachyte hold rocks richer in
alkalis at the same silica; the U fields, phonolite and foidite hold rocks so rich in
alkalis, and so poor in silica, that they crystallise feldspathoids. A few fields are
subdivided by a second criterion: a trachybasalt is called hawaiite if $\ce{Na2O} - 2
\ge \ce{K2O}$, potassic trachybasalt otherwise; a tephrite with more than
\qty{10}{\percent} olivine in its norm (below) is a basanite; a rhyolite or a trachyte whose
molecular $(\ce{Na2O + K2O})/\ce{Al2O3}$ exceeds 1 is called peralkaline. The scheme
asks for analyses with less than \qty{2}{\percent} water and \qty{0.5}{\percent}
carbon dioxide: alteration adds water and moves the alkalis, and a \omterm{def:b2:magmatic-differentiation:cumulate}{cumulate} is not a
liquid.
% ledger: bgs:tas, bgs:tas-use

\begin{omfigure}
\centering
\begin{tikzpicture}
\begin{axis}[omrecord, width=0.86\linewidth, height=8.6cm,
  xmin=37, xmax=80, ymin=0, ymax=16, xtick={37,41,...,77}, ytick={0,2,...,16},
  xlabel={\ce{SiO2} (\unit{\percent})}, ylabel={\ce{Na2O + K2O} (\unit{\percent})},
  legend style={font=\scriptsize, at={(0.98,0.98)}, anchor=north east},
  tick label style={font=\scriptsize}, label style={font=\small}]
\addplot[only marks, mark=*, mark size=0.55pt, blue!70!black, opacity=0.6] table[x=SiO2, y=alk] {figdata/out/bachelor-2-series-cath-th.dat};
\addlegendentry{Iceland}
\addplot[only marks, mark=*, mark size=0.55pt, red!75!black, opacity=0.6] table[x=SiO2, y=alk] {figdata/out/bachelor-2-series-cath-ca.dat};
\addlegendentry{Cascades}
\foreach \f in {Pc,B,O1,O2,O3,R,S1,S2,S3,T,U1,U2,U3,Ph,F}
  \addplot[black, thin, forget plot] table[x=SiO2, y=alk] {figdata/out/bachelor-2-series-tas-\f.dat};
\node[font=\tiny] at (axis cs:43,1.6) {Pc};
\node[font=\tiny, fill=white, inner sep=0.5pt] at (axis cs:50.5,0.8) {B};
\node[font=\tiny] at (axis cs:54.5,2.2) {O1};
\node[font=\tiny] at (axis cs:60,2.2) {O2};
\node[font=\tiny] at (axis cs:66.5,2.2) {O3};
\node[font=\tiny] at (axis cs:74.5,11) {R};
\node[font=\tiny, fill=white, inner sep=0.5pt] at (axis cs:49,5.8) {S1};
\node[font=\tiny] at (axis cs:52.5,7.4) {S2};
\node[font=\tiny] at (axis cs:57.5,8.9) {S3};
\node[font=\tiny] at (axis cs:64.5,12) {T};
\node[font=\tiny] at (axis cs:44,6.6) {U1};
\node[font=\tiny] at (axis cs:48.7,9.4) {U2};
\node[font=\tiny] at (axis cs:52.5,11.8) {U3};
\node[font=\tiny] at (axis cs:57,14.6) {Ph};
\node[font=\tiny] at (axis cs:39,10) {F};
\end{axis}
\end{tikzpicture}
\par\smallskip{\small The TAS diagram (after Le Bas and others, 1986, as given by the
British Geological Survey) with 456 lavas and intrusive rocks from Iceland (blue) and
580 from the Cascades (red), oxides as analysed. Pc picrobasalt, B basalt, O1 basaltic
andesite, O2 andesite, O3 dacite, R rhyolite; S1 trachybasalt, S2 basaltic
trachyandesite, S3 trachyandesite, T trachyte or trachydacite; U1 tephrite or
basanite, U2 phonotephrite, U3 tephriphonolite, Ph phonolite, F foidite.}
\end{omfigure}
% ledger: bgs:tas, geoplotr:tas, data:b4-cath


\begin{definition}[Silica saturation]\label{def:b2:magmatic-series:saturation}
A \omterm{def:b2:magma-emplacement:magma}{magma} or rock is \emph{silica-oversaturated}\index{silica saturation} if it has
more silica than its other oxides can combine with in silicate minerals, so that
quartz (or another silica mineral) appears; \emph{silica-undersaturated} if it has too
little to turn all its alkalis into feldspar, so that a feldspathoid (nepheline,
leucite) appears instead, never together with quartz; and \emph{silica-saturated}
between the two.
\end{definition}

For a lava whose minerals cannot be counted, saturation is judged on a \emph{norm}: a
standard calculation that distributes the oxides of the analysis among ideal minerals
in a fixed order, alkalis into feldspar first, and reports the excess silica as
quartz or the shortfall as feldspathoid. The usual one, the CIPW norm, separates the
two kinds of basalt: a basalt with nepheline in its norm is an alkali basalt; one
without is a subalkali basalt, and among these the tholeiitic basalts carry
hypersthene (orthopyroxene) in the norm and, often, low-calcium pyroxene in the mode.
\Cref{exo:b2:magmatic-series:11} runs a simplified version of it.
% ledger: bgs:tas-use, bgs:series-terms

\section{Magmatic series}\label{sec:b2:magmatic-series:series}

Lavas from one volcano, or one province, are seldom alike, but they are related: they
derive from similar parent \omterm{def:b2:magma-emplacement:magma}{magmas} by the differentiation of
\cref{ch:b2:magmatic-differentiation}, and they trace together a line of descent.

\begin{definition}[Magmatic series]\label{def:b2:magmatic-series:series}
A \emph{magmatic series}\index{magmatic series} is a suite of igneous rocks of one
magmatic cycle or province related by differentiation from similar parent \omterm{def:b2:magma-emplacement:magma}{magmas} and
sharing characteristic chemical and mineralogical features. Three are commonly
distinguished. The \emph{tholeiitic series}\index{tholeiitic series} starts from
basalts saturated or slightly oversaturated in silica, with hypersthene in the norm,
and shows a marked enrichment in iron at the start of its differentiation. The
\emph{calc-alkaline series}\index{calc-alkaline series} runs from basalt through
andesite to dacite and rhyolite, with little or no enrichment in iron, and with
hydrous minerals (hornblende, biotite) in its evolved members. The \emph{alkaline
series}\index{alkaline series} starts from nepheline-normative basalts and evolves
through the S and U fields of the TAS diagram towards trachyte and phonolite.
\end{definition}
% ledger: bgs:series-terms, bgs:tas-use

The names describe suites, not single rocks: an isolated basalt can belong to any of
them, and the same names have been used with different meanings, so a petrologist
must say what she means by them. The distinction between the tholeiitic and
\omterm{def:b2:magmatic-series:series}{calc-alkaline series} rests on iron, and the diagram that shows it best plots three
groups of oxides.
% ledger: bgs:series-terms

\begin{definition}[AFM diagram]\label{def:b2:magmatic-series:afm}
An \emph{AFM diagram}\index{AFM diagram} is a triangular diagram whose corners are
$A = \ce{Na2O + K2O}$, $F = \ce{FeO_T}$ (all iron as \ce{FeO}) and $M = \ce{MgO}$, in
mass per cent recalculated so that $A + F + M = 100$. A line drawn on it separates the
tholeiitic field, on the side of F, from the calc-alkaline field.
\end{definition}
% ledger: vp21:series, geoplotr:afm

A differentiating basalt loses magnesium faster than anything else, since olivine and
pyroxene are richer in magnesium than the liquid, and keeps its alkalis, which are
incompatible: its line of descent moves away from M, towards A. What happens to
iron decides the series. In a tholeiitic \omterm{def:b2:magma-emplacement:magma}{magma} iron leaves only with the
magnesium-rich ferrous silicates, more slowly than magnesium, so that the ratio
$F/M$ rises and the trend first climbs towards F before turning to A: the
\emph{Fenner trend}. In a calc-alkaline \omterm{def:b2:magma-emplacement:magma}{magma}, more oxidised, iron also leaves in
magnetite, which holds ferric iron, and $F/M$ stays nearly constant as the liquid
moves straight towards A: the \emph{Bowen trend}. The figure below shows both with
real rocks, Iceland's tholeiites and the calc-alkaline lavas of the Cascades. The
dividing line, drawn by eye in 1971, separates them well where they are basaltic; at
high alkalis the two trends converge and the line says little.
% ledger: vp21:series, geoplotr:afm, data:b4-cath

\begin{omfigure}
\centering
\begin{tikzpicture}
\begin{axis}[name=afm, width=0.5\linewidth, height=0.5\linewidth, axis equal image,
  hide axis, xmin=-0.04, xmax=1.04, ymin=-0.06, ymax=0.92, clip=false]
\addplot[only marks, mark=*, mark size=0.5pt, blue!70!black, opacity=0.6] table[x=afm_x, y=afm_y] {figdata/out/bachelor-2-series-cath-th.dat};
\addplot[only marks, mark=*, mark size=0.5pt, red!75!black, opacity=0.6] table[x=afm_x, y=afm_y] {figdata/out/bachelor-2-series-cath-ca.dat};
\addplot[very thick, black] table[x=x, y=y] {figdata/out/bachelor-2-series-afm-line.dat};
\draw[thick] (axis cs:0,0) -- (axis cs:1,0) -- (axis cs:0.5,0.8660254) -- cycle;
\node[anchor=south] at (axis cs:0.5,0.866) {F};
\node[anchor=north east] at (axis cs:0,0) {A};
\node[anchor=north west] at (axis cs:1,0) {M};
\node[font=\scriptsize, fill=white, inner sep=1pt] at (axis cs:0.52,0.62) {tholeiitic};
\node[font=\scriptsize, fill=white, inner sep=1pt] at (axis cs:0.33,0.12) {calc-alkaline};
\end{axis}
\begin{axis}[omrecord, at={(afm.east)}, anchor=west, xshift=1.3cm,
  width=0.46\linewidth, height=0.46\linewidth,
  xmin=-2.5, xmax=2.2, ymin=-4.5, ymax=1.5,
  xlabel={$\ln(A/F)$}, ylabel={$\ln(M/F)$},
  legend style={font=\scriptsize, at={(0.98,0.98)}, anchor=north east},
  tick label style={font=\scriptsize}, label style={font=\small}]
\addplot[only marks, mark=*, mark size=0.5pt, blue!70!black, opacity=0.6] table[x=lnAF, y=lnMF] {figdata/out/bachelor-2-series-lr-th.dat};
\addlegendentry{Iceland}
\addplot[only marks, mark=*, mark size=0.5pt, red!75!black, opacity=0.6] table[x=lnAF, y=lnMF] {figdata/out/bachelor-2-series-lr-ca.dat};
\addlegendentry{Cascades}
\addplot[thick, blue!50!black, forget plot] table[x=x, y=y_th] {figdata/out/bachelor-2-series-lr-fits.dat};
\addplot[thick, red!50!black, forget plot] table[x=x, y=y_ca] {figdata/out/bachelor-2-series-lr-fits.dat};
\end{axis}
\end{tikzpicture}
\par\smallskip{\small Left: the \omterm{def:b2:magmatic-series:afm}{AFM diagram} of the Icelandic (blue) and Cascades (red)
rocks of the TAS figure, with the dividing line of Irvine and Baragar (1971). Right:
the same rocks in log-ratio coordinates, where each series becomes a nearly straight
band; the lines are orthogonal regressions, of slopes $-1.10$ (Iceland) and $-0.72$
(Cascades), meeting at $(-3.84, 2.12)$ as Vermeesch and Pease (2021) found.}
\end{omfigure}
% ledger: geoplotr:afm, data:b4-cath, vp21:model

Of the Icelandic rocks within the range of the dividing line, 351 fall on the
tholeiitic side and 23 on the other; of the Cascades rocks, 417 on the calc-alkaline
side and 57 on the other. A curve is a clumsy boundary, and a simple model explains
the shape of the trends. Suppose that a \omterm{def:b2:magma-emplacement:magma}{magma} loses its alkalis, iron and magnesium
at rates proportional to the amounts it holds, $\mathrm dX/\mathrm dt = -\lambda_X X$.

\begin{proposition}[Log-ratio trends]\label{prop:b2:magmatic-series:logratio}
If $A$, $F$ and $M$ decay exponentially with constants $\lambda_A$, $\lambda_F$ and
$\lambda_M$, the \omterm{def:b2:magma-emplacement:magma}{magma} moves along a straight line in the plane of $\ln(A/F)$ and
$\ln(M/F)$, of slope $(\lambda_F - \lambda_M)/(\lambda_F - \lambda_A)$, while it moves
along a curve on the AFM triangle.
\end{proposition}

\begin{proof}
With $X = X_0 e^{-\lambda_X t}$, $\ln(A/F) = \ln(A_0/F_0) + (\lambda_F - \lambda_A)t$
and $\ln(M/F) = \ln(M_0/F_0) + (\lambda_F - \lambda_M)t$: both are linear in $t$, so
one is linear in the other, with the slope stated. The ratios $A/F$ and $M/F$ fix the
point on the triangle, whose coordinates $A/(A + F + M)$, \dots\ are not linear in $t$.
\end{proof}

The steeper slope of the Icelandic rocks means a smaller $\lambda_F$: iron leaves
the tholeiitic \omterm{def:b2:magma-emplacement:magma}{magma} more slowly (\cref{exo:b2:magmatic-series:10}). Vermeesch and Pease showed that the two
bands of the figure above meet at one point, which they read as a hypothetical common
source.
% ledger: vp21:model

The tholeiitic and \omterm{def:b2:magmatic-series:series}{calc-alkaline series} differ in their potash as well. The
\ce{K2O}--\ce{SiO2} diagram divides basalts, andesites, dacites and rhyolites with
more than \qty{48}{\percent} silica into low-K, medium-K and high-K types by two
straight lines (the figure below). Most of both suites here are medium-K: the
diagram says how much potash a series carries, not which series it is.
% ledger: bgs:k2o-sio2, data:b4-cath

\begin{omfigure}
\centering
\begin{tikzpicture}
\begin{axis}[omrecord, width=0.78\linewidth, height=6.2cm,
  xmin=45, xmax=79, ymin=0, ymax=5,
  xlabel={\ce{SiO2} (\unit{\percent})}, ylabel={\ce{K2O} (\unit{\percent})},
  legend style={font=\scriptsize, at={(0.02,0.98)}, anchor=north west},
  tick label style={font=\scriptsize}, label style={font=\small}]
\addplot[only marks, mark=*, mark size=0.55pt, blue!70!black, opacity=0.6] table[x=SiO2, y=K2O] {figdata/out/bachelor-2-series-cath-th.dat};
\addlegendentry{Iceland}
\addplot[only marks, mark=*, mark size=0.55pt, red!75!black, opacity=0.6] table[x=SiO2, y=K2O] {figdata/out/bachelor-2-series-cath-ca.dat};
\addlegendentry{Cascades}
\addplot[thick, black, forget plot] table[x=x, y=y_low] {figdata/out/bachelor-2-series-k2o-lines.dat};
\addplot[thick, black, forget plot] table[x=x, y=y_high] {figdata/out/bachelor-2-series-k2o-lines.dat};
\draw[thick] (axis cs:48,0) -- (axis cs:48,0.3);
\node[font=\scriptsize, fill=white, inner sep=1pt] at (axis cs:60,3.7) {high-K};
\node[font=\scriptsize, fill=white, inner sep=1pt] at (axis cs:76.5,2.2) {medium-K};
\node[font=\scriptsize, fill=white, inner sep=1pt] at (axis cs:60,0.35) {low-K};
\end{axis}
\end{tikzpicture}
\par\smallskip{\small Potash against silica for the Icelandic and Cascades rocks, with
the lines through (48, 0.3)--(68, 1.2) and (48, 1.2)--(68, 3.1) that divide low-K,
medium-K and high-K types (after Peccerillo and Taylor, 1976, as adopted by the
international scheme), extended here beyond 68.}
\end{omfigure}
% ledger: bgs:k2o-sio2, data:b4-cath

Why does oxidation differ? Calc-alkaline \omterm{def:b2:magma-emplacement:magma}{magmas} form above subduction zones, where
the dehydrating slab supplies oxidised species to the mantle wedge; the \omterm{def:b2:magma-emplacement:magma}{magmas}
crystallise magnetite, and their evolved members hornblende or biotite. Tholeiitic
\omterm{def:b2:magma-emplacement:magma}{magmas} form at ridges, where the mantle melts by decompression
(\cref{ch:b2:mantle-melting}), and are more reduced.
% ledger: vp21:series

\section{Series and settings}\label{sec:b2:magmatic-series:settings}

Each setting of plate tectonics produces its own \omterm{def:b2:magma-emplacement:magma}{magmas}, because each melts the
mantle in its own way: by decompression beneath ridges and hotspots, by adding water
above subduction zones (\cref{ch:b2:mantle-melting}).

\emph{Mid-ocean ridges} erupt tholeiitic basalts, the MORB of
\cref{ch:b2:elements-in-the-earth}, poor in \omterm{def:b2:elements-in-the-earth:partition}{incompatible elements} and with the
depleted isotopic signature of their source. Iceland, a ridge over a hotspot, adds
its own tholeiites to that list, with a full suite to rhyolite.
% ledger: vp21:series

\emph{Ocean islands} grow in stages. A Hawaiian volcano erupts alkalic lavas while
young, then builds its shield of tholeiitic basalt, an estimated 80 to
\qty{95}{\percent} of its final volume, within its first 1 to 2 million years; as it
drifts away from the hotspot its lavas become alkalic again, ending in a long tail of
rare alkalic eruptions. Less melting gives more alkaline \omterm{def:b2:magma-emplacement:magma}{magmas}: the alkalis, being
incompatible, are diluted less in a small \omterm{def:b2:mantle-melting:fraction}{melt fraction}. Isotopically, ocean-island
basalts sample the enriched reservoirs of \cref{ch:b2:isotope-geochemistry}.
% ledger: usgs14:hawaii-stages

\emph{Volcanic arcs} erupt \omterm{def:b2:magmatic-series:series}{calc-alkaline series}. Their \omterm{def:b2:elements-in-the-earth:trace}{trace elements} carry a
distinctive mark (the figure below): normalised to the \omterm{def:b2:elements-in-the-earth:bse}{primitive mantle}, the arc
basalts are enriched in Rb, Ba, Th, U, Pb and Sr, elements that the slab adds to the
wedge, and depleted in niobium and tantalum relative to their neighbours. The median
arc basalt of the data set has $\mathrm{Nb}/\mathrm{La} = 0.49$, against 0.97 for MORB
and 1.23 for ocean-island basalts.
% ledger: data:b4-basalt-settings

\begin{omfigure}
\centering
\begin{tikzpicture}
\begin{axis}[omrecord, width=0.86\linewidth, height=6.4cm, ymode=log,
  xmin=0.5, xmax=22.5, ymin=1, ymax=300,
  xtick={1,...,22}, xticklabels={Rb,Ba,Th,U,Nb,Ta,La,Ce,Pb,Sr,Nd,Zr,Hf,Sm,Eu,Ti,Gd,Dy,Y,Er,Yb,Lu},
  ylabel={sample / primitive mantle},
  legend style={font=\scriptsize, at={(0.98,0.98)}, anchor=north east},
  tick label style={font=\scriptsize}, label style={font=\small}]
\addplot[thick, mark=square*, mark size=1.4pt, orange!85!black] table[x=idx, y=oib] {figdata/out/bachelor-2-series-spider.dat};
\addlegendentry{ocean-island basalt}
\addplot[thick, mark=*, mark size=1.4pt, blue!70!black] table[x=idx, y=morb] {figdata/out/bachelor-2-series-spider.dat};
\addlegendentry{MORB}
\addplot[thick, mark=triangle*, mark size=1.7pt, red!75!black] table[x=idx, y=iab] {figdata/out/bachelor-2-series-spider.dat};
\addlegendentry{island-arc basalt}
\end{axis}
\end{tikzpicture}
\par\smallskip{\small Primitive-mantle-normalised median compositions of 241 MORB, 259
ocean-island basalts and 256 island-arc basalts (the compilation of Vermeesch, 2006;
\omterm{def:b2:elements-in-the-earth:bse}{primitive mantle} of McDonough and Sun, 1995). The arc basalts show the trough at Nb
and Ta and peaks at Ba, Pb and Sr; the ocean-island basalts are enriched in all the
most \omterm{def:b2:elements-in-the-earth:partition}{incompatible elements}.}
\end{omfigure}
% ledger: data:b4-basalt-settings, ms95:pm-trace

Within continents, basaltic \omterm{def:b2:magma-emplacement:magma}{magma} meets a thick crust that it can heat and melt,
and the result can be a volcanism with two modes.

\begin{definition}[Bimodal volcanism]\label{def:b2:magmatic-series:bimodal}
\emph{Bimodal volcanism}\index{bimodal volcanism} is the eruption, in one province,
of two groups of \omterm{def:b2:magma-emplacement:magma}{magmas}, basaltic and rhyolitic, with few intermediate compositions
between them.
\end{definition}

The volcanism along the track of the Yellowstone hotspot is bimodal: on the Eastern
Snake River Plain, older rhyolites with some \qty{75}{\percent} silica are covered by
younger basalts with some \qty{50}{\percent}. The gap between them is the clue: the
differentiation of a basalt alone would have left intermediate lavas behind, and a gap
points to a second source for the rhyolites, such as the crust that the basalt heats.
% ledger: yvo23:bimodal

\emph{Large igneous provinces}, met in the Year 1 volume, erupt basalt in volumes
no other setting matches; their series is read with the same diagrams.

\section{Discrimination diagrams}\label{sec:b2:magmatic-series:discrimination}

The patterns of the previous section suggest reading a setting from an analysis
directly, and petrologists have drawn many diagrams to do so.

\begin{definition}[Discrimination diagram]\label{def:b2:magmatic-series:discrimination}
A \emph{discrimination diagram}\index{discrimination diagram} is a diagram of
chemical variables, usually immobile \omterm{def:b2:elements-in-the-earth:trace}{trace elements}, divided into fields drawn
around the compositions of rocks from known settings, used to assign a setting to a
rock of unknown origin.
\end{definition}

The classic is the Ti--Zr--Y triangle of Pearce and Cann (1973), drawn for basalts. It
uses elements little moved by alteration, so that it can be applied to old,
metamorphosed lavas whose setting is lost. The figure below tests it on the 538 basalts of the
compilation that have all three elements. Most ocean-island basalts fall in the
within-plate field D, 152 of 173, and most MORB in field B, 142 of 164; but field B is
shared with arc basalts, and 101 of the 201 arc basalts fall there too.
% ledger: geoplotr:tizry, data:b4-basalt-settings

\begin{omfigure}
\centering
\begin{tikzpicture}
\begin{axis}[width=0.62\linewidth, height=0.55\linewidth, axis equal image,
  hide axis, xmin=-0.04, xmax=1.04, ymin=-0.06, ymax=0.92, clip=false,
  legend style={font=\scriptsize, at={(1.0,0.95)}, anchor=north west}]
\addplot[only marks, mark=*, mark size=0.7pt, blue!70!black, opacity=0.7] table[x=x, y=y] {figdata/out/bachelor-2-series-tizry-morb.dat};
\addlegendentry{MORB}
\addplot[only marks, mark=square*, mark size=0.7pt, orange!85!black, opacity=0.7] table[x=x, y=y] {figdata/out/bachelor-2-series-tizry-oib.dat};
\addlegendentry{ocean-island basalt}
\addplot[only marks, mark=triangle*, mark size=0.9pt, red!75!black, opacity=0.7] table[x=x, y=y] {figdata/out/bachelor-2-series-tizry-iab.dat};
\addlegendentry{island-arc basalt}
\foreach \f in {A,B,C,D}
  \addplot[thick, black, forget plot] table[x=x, y=y] {figdata/out/bachelor-2-series-tizry-\f.dat};
\draw[thick] (axis cs:0,0) -- (axis cs:1,0) -- (axis cs:0.5,0.8660254) -- cycle;
\node[anchor=south] at (axis cs:0.5,0.866) {Ti/100};
\node[anchor=north east] at (axis cs:0,0) {Zr};
\node[anchor=north west] at (axis cs:1,0) {Y $\times$ 3};
\draw[thin] (axis cs:0.33,0.36) -- (axis cs:0.22,0.46) node[anchor=south east, font=\small\bfseries, inner sep=1pt] {D};
\draw[thin] (axis cs:0.36,0.12) -- (axis cs:0.24,0.06) node[anchor=north east, font=\small\bfseries, inner sep=1pt] {C};
\draw[thin] (axis cs:0.56,0.37) -- (axis cs:0.76,0.50) node[anchor=west, font=\small\bfseries, inner sep=1pt] {A};
\draw[thin] (axis cs:0.53,0.21) -- (axis cs:0.72,0.13) node[anchor=north west, font=\small\bfseries, inner sep=1pt] {B};
\end{axis}
\end{tikzpicture}
\par\smallskip{\small The Ti--Zr--Y diagram of Pearce and Cann (1973) with the basalts
of known setting. A island-arc tholeiites; B MORB, island-arc tholeiites and
calc-alkali basalts; C calc-alkali basalts; D within-plate basalts. Ti from
\ce{TiO2}; all three in ppm.}
\end{omfigure}
% ledger: geoplotr:tizry, data:b4-basalt-settings

\begin{method}[Using a discrimination diagram]\label{met:b2:magmatic-series:discrimination}
(1) Check that the rock is of the kind the diagram was built for, usually a basalt, and
not a \omterm{def:b2:magmatic-differentiation:cumulate}{cumulate} or an evolved lava, whose \omterm{def:b2:elements-in-the-earth:trace}{trace elements} reflect crystallisation as much
as the source. (2) Prefer elements immobile in alteration (Ti, Zr, Y, Nb, the heavy
rare earths) for altered or metamorphosed rocks. (3) Plot a suite of samples, not one:
a single point near a boundary decides nothing. (4) Read the result as a likelihood:
the fields were drawn around most, not all, of their reference samples, and some fields
are shared. (5) Confirm with independent evidence: other diagrams, isotopes, the
geological context.
\end{method}

Step (4) can be made quantitative (\cref{exo:b2:magmatic-series:8}). If a basalt falls in field B, and the
three settings were equally common, the chance that it is a MORB is only about 0.62.
Statistical classifiers trained on such compilations inherit the same limit: a basalt
carries the memory of its source and its melting, and two settings can make similar
\omterm{def:b2:magma-emplacement:magma}{magmas}.
% ledger: data:b4-basalt-settings

\begin{inthelab}[Checking a whole-rock analysis]
Before plotting a rock on any of these diagrams, petrologists check the analysis. The
major oxides, measured for instance by X-ray fluorescence, should add up close to 100
with the volatiles: water and carbon dioxide, measured directly or together as the
loss on ignition, the mass a powder loses on strong heating. A large loss on
ignition marks an altered rock, whose alkalis may have moved; the scheme of the TAS
diagram rejects analyses with more than \qty{2}{\percent} water or
\qty{0.5}{\percent} carbon dioxide. Analyses are then recalculated to 100 free of
volatiles, iron is expressed as total \ce{FeO} or split between ferrous and ferric,
and reference materials analysed in the same run are compared with their certified
values (\cref{ch:b2:field-lab-methods-2}).
\end{inthelab}
% ledger: bgs:tas-use

\section{Exercises}

\begin{exercise}[$\star$]\label{exo:b2:magmatic-series:1}
Name the following volcanic rocks on the TAS diagram (\ce{SiO2}, \ce{Na2O}, \ce{K2O}
in \unit{\percent}, already volatile-free): (a) 49.60, 2.22, 0.51 (the Hawaiian
reference basalt BHVO-2); (b) 58.5, 3.9, 1.6; (c) 72.0, 4.0, 4.2; (d) 47.0, 4.5, 2.5.
For (d), name the two possible rocks and say what decides between them.
\end{exercise}
% ledger: usgs:bhvo2-major, bgs:tas

\begin{exercise}[$\star$]\label{exo:b2:magmatic-series:2}
A lava gives \ce{SiO2} 51.60, \ce{TiO2} 1.20, \ce{Al2O3} 16.10, \ce{FeO_T} 9.80,
\ce{MgO} 6.60, \ce{CaO} 8.20, \ce{Na2O} 3.50, \ce{K2O} 1.40, \ce{H2O} 1.40 and
\ce{CO2} \qty{0.20}{\percent}. (a) Is the analysis acceptable for the TAS diagram?
(b) Recalculate it free of volatiles. (c) Name the rock before and after the
recalculation. (d) Is it low-, medium- or high-K?
\end{exercise}
% ledger: bgs:tas-use, bgs:k2o-sio2

\begin{exercise}[$\star$]\label{exo:b2:magmatic-series:3}
Name on the \omterm{def:b2:magmatic-series:qapf}{QAPF diagram}: (a) a coarse-grained rock with \qty{25}{\percent} quartz,
\qty{30}{\percent} alkali feldspar, \qty{35}{\percent} plagioclase, \qty{8}{\percent}
biotite and \qty{2}{\percent} hornblende; (b) one with \qty{5}{\percent} quartz,
\qty{5}{\percent} alkali feldspar, \qty{55}{\percent} plagioclase and
\qty{35}{\percent} clinopyroxene.
\end{exercise}
% ledger: bgs:qapf

\begin{exercise}[$\star$]\label{exo:b2:magmatic-series:4}
Classify as low-, medium- or high-K, with the lines of the \ce{K2O}--\ce{SiO2} figure:
(a) 50.0, 0.40; (b) 56.0, 1.50; (c) 62.0, 2.90; (d) 66.0, 1.00 (\ce{SiO2},
\ce{K2O}, \unit{\percent}).
\end{exercise}
% ledger: bgs:k2o-sio2

\begin{exercise}[$\star\star$]\label{exo:b2:magmatic-series:5}
(a) Compute $A$, $F$ and $M$ for a basalt with \ce{FeO_T} 14.5, \ce{MgO} 4.5 and
\ce{Na2O + K2O} \qty{3.2}{\percent}, and for an andesite with 6.1, 3.8 and 5.05.
(b) The dividing line passes through $(F, A, M) = (54.5, 23.0, 22.5)$, $(52.5, 30.0,
17.5)$ and $(49.6, 35.0, 15.4)$, the first being its most iron-rich point; the
tholeiitic side is the side of F. On which side does each rock fall? (c) Which
mineral, crystallising early, keeps the andesite where it is?
\end{exercise}
% ledger: geoplotr:afm, vp21:series

\begin{exercise}[$\star\star$]\label{exo:b2:magmatic-series:6}
The median MORB of the compilation has \qty{4}{ppm} Nb and \qty{4.12}{ppm} La, the
median ocean-island basalt 38 and 30.8, the median arc basalt 3 and 6.1. With the
\omterm{def:b2:elements-in-the-earth:bse}{primitive mantle} values Nb 0.658 and La \qty{0.648}{ppm}, compute the normalised values
and $\mathrm{Nb_N}/\mathrm{La_N}$ for each. Which has a niobium anomaly, and why?
\end{exercise}
% ledger: data:b4-basalt-settings, ms95:pm-trace

\begin{exercise}[$\star\star$]\label{exo:b2:magmatic-series:7}
Ti is $0.5993$ of the mass of \ce{TiO2}. Compute the Ti--Zr--Y plot coordinates
(Ti/100, Zr, $3\mathrm Y$, recast to 100) of three basalts: (a) \ce{TiO2}
\qty{1.6}{\percent}, Zr 110, Y \qty{32}{ppm}; (b) 2.8, 220, 29; (c) 0.8, 50, 20. Locate
them on the figure and say what each placement allows you to conclude.
\end{exercise}
% ledger: geoplotr:tizry

\begin{exercise}[$\star\star$]\label{exo:b2:magmatic-series:8}
In the Ti--Zr--Y test of this chapter, field B holds 142 of 164 MORB, 4 of 173
ocean-island basalts and 101 of 201 arc basalts. (a) What fraction of each group falls
in B? (b) If the three settings were equally common, what is the probability that a
basalt plotting in B is a MORB? (c) What does field B alone allow you to say?
\end{exercise}
% ledger: data:b4-basalt-settings

\begin{exercise}[$\star\star$]\label{exo:b2:magmatic-series:9}
A rhyolite has \ce{Na2O} 5.5, \ce{K2O} 4.5 and \ce{Al2O3} \qty{10.5}{\percent} ($M$ =
61.98, 94.20 and \qty{101.96}{g/mol}). Compute its molecular $(\ce{Na2O +
K2O})/\ce{Al2O3}$. (a) Is it peralkaline? (b) Alkali feldspar holds one
alkali oxide per \ce{Al2O3}: what fraction of the alkalis cannot enter it?
\end{exercise}
% ledger: bgs:tas-use

\begin{exercise}[$\star\star\star$]\label{exo:b2:magmatic-series:10}
In the model of the log-ratio proposition, take $\lambda_A = 0.4$ and $\lambda_M =
0.8$. (a) Compute the slope of the trend for $\lambda_F = 0.5$ and for $\lambda_F =
0.7$. (b) The Icelandic rocks have a slope of $-1.10$ and those of the Cascades
$-0.72$: what $\lambda_F$ does each imply? (c) Which series loses its iron faster,
and which mineral explains it? (d) What happens to $F/M$ along each trend?
\end{exercise}
% ledger: vp21:model

\begin{exercise}[$\star\star\star$]\label{exo:b2:magmatic-series:11}
A simplified norm. Work in moles of oxide. Put all \ce{K2O} into orthoclase
(\ce{K2O.Al2O3.6SiO2}), all \ce{Na2O} into albite (\ce{Na2O.Al2O3.6SiO2}), the
remaining \ce{Al2O3} with as much \ce{CaO} into anorthite (\ce{CaO.Al2O3.2SiO2}), the
remaining \ce{CaO} with as much $(\ce{Mg,Fe})\ce{O}$ into diopside
(\ce{CaO.(Mg,Fe)O.2SiO2}) and the rest of $(\ce{Mg,Fe})\ce{O}$ into hypersthene
(\ce{(Mg,Fe)O.SiO2}). If silica is short, turn hypersthene into olivine
(\ce{2(Mg,Fe)O.SiO2}); if it is still short, turn albite into nepheline
(\ce{Na2O.Al2O3.2SiO2}). Apply it to (a) BHVO-2: \ce{SiO2} 49.60, \ce{Al2O3} 13.44,
\ce{FeO_T} 11.15, \ce{MgO} 7.257, \ce{CaO} 11.40, \ce{Na2O} 2.219, \ce{K2O}
\qty{0.513}{\percent}; (b) a basalt with 45.5, 15.0, 11.5, 8.5, 10.0, 3.8 and 1.5 (molar
masses \ce{SiO2} 60.08, \ce{Al2O3} 101.96, \ce{FeO} 71.84, \ce{MgO} 40.30, \ce{CaO}
56.08, \ce{Na2O} 61.98, \ce{K2O} \qty{94.20}{g/mol}). Is each silica-saturated? To
which series does each basalt belong?
\end{exercise}
% ledger: usgs:bhvo2-major, bgs:series-terms

\begin{exercise}[$\star\star\star$]\label{exo:b2:magmatic-series:12}
In the compilation, $\mathrm{Nb}/\mathrm{La} < 0.7$ holds for 116 of 141 arc basalts,
21 of 101 MORB and 3 of 137 ocean-island basalts analysed for both. (a) Compute the
three fractions. (b) With equal proportions of the three settings, what is the
probability that a basalt with $\mathrm{Nb}/\mathrm{La} < 0.7$ is an arc basalt? (c)
Redo it for a region where \qty{60}{\percent} of basalts are MORB, \qty{30}{\percent}
ocean-island basalts and \qty{10}{\percent} arc basalts. (d) What lesson does (c)
hold for the use of \omterm{def:b2:magmatic-series:discrimination}{discrimination diagrams}?
\end{exercise}
% ledger: data:b4-basalt-settings

\section{Problem: One volcano, three lavas}

\begin{problem}[{Weekend problem --- from three analyses to the setting of a volcano and
the crystals it left behind}]
\label{pb:b2:magmatic-series:1}
Three lavas of an (invented) stratovolcano, from its oldest flank flows to its youngest
dome, give (\unit{\percent}):

\smallskip\noindent{\small
\begin{tabular}{@{}lrrrrrrrrrr@{}}
 & \ce{SiO2} & \ce{TiO2} & \ce{Al2O3} & \ce{FeO_T} & \ce{MgO} & \ce{CaO} & \ce{Na2O} & \ce{K2O} & \ce{H2O} & \ce{CO2}\\
lava 1 & 50.4 & 0.95 & 17.6 & 8.6 & 7.9 & 10.2 & 2.9 & 0.62 & 0.70 & 0.10\\
lava 2 & 58.6 & 0.80 & 17.4 & 6.1 & 3.8 & 6.9 & 3.9 & 1.15 & 1.10 & 0.05\\
lava 3 & 65.3 & 0.55 & 16.2 & 4.0 & 1.6 & 4.3 & 4.4 & 1.85 & 1.20 & 0.05\\
\end{tabular}}

\smallskip\noindent and in ppm: Rb 10, 20, 32; Th 1.1, 2.1, 3.4; Nb 2.0, 3.6, 5.6;
La 7, 13, 20; Sr 520, 480, 380; Zr 70, 125, 190; Y 18, 22, 24. \omterm{def:b2:elements-in-the-earth:bse}{Primitive mantle}: Nb
0.658, La 0.648, Th \qty{0.0795}{ppm}.

\noindent\textbf{Part I --- Naming the lavas.}
\begin{enumerate}
\item Add up each analysis. Are the three acceptable for the TAS diagram?
\item Recalculate lava 1 free of volatiles: give the factor, \ce{SiO2} and
\ce{Na2O + K2O}.
\item Do the same for lavas 2 and 3.
\item Name the three lavas on the TAS diagram.
\item Would the names change without the recalculation?
\item A coarse-grained rock of the composition of lava 2 crystallised in the volcano's
roots. Which scheme names it, and in which field would you expect it?
\end{enumerate}

\noindent\textbf{Part II --- Which series?}
\begin{enumerate}[resume]
\item Classify each lava as low-, medium- or high-K.
\item Compute $A$, $F$, $M$ for each (volatile-free).
\item The dividing line passes through $(F, A, M) = (42.8, 17.1, 40.1)$, $(46.9, 18.1,
35.0)$, $(52.5, 30.0, 17.5)$, $(49.6, 35.0, 15.4)$, $(40.8, 50.0, 9.2)$ and $(37.9,
55.0, 7.1)$, with the tholeiitic field on the side of F. Interpolating $F$ linearly in
$A$ between neighbouring points, on which side does each lava fall?
\item Compute \ce{FeO_T}/\ce{MgO} for each lava. Is there an iron enrichment?
\item To which series does the volcano belong? Which mineral and which conditions
explain it?
\end{enumerate}

\noindent\textbf{Part III --- Which setting?}
\begin{enumerate}[resume]
\item Normalise Nb, La and Th of lava 1 to the \omterm{def:b2:elements-in-the-earth:bse}{primitive mantle}.
\item Compute $\mathrm{Nb_N}/\mathrm{La_N}$ for lava 1 and compare it with the medians
of \cref{exo:b2:magmatic-series:6}.
\item Compute $\mathrm{Th_N}/\mathrm{Nb_N}$. What carried the thorium?
\item Compute the Ti--Zr--Y coordinates of lava 1 (\ce{TiO2} of the volatile-free
analysis). What does its field say? Why not plot lavas 2 and 3?
\item Conclude: where did this volcano grow?
\end{enumerate}

\noindent\textbf{Part IV --- How much crystallised?}
\begin{enumerate}[resume]
\item Potash enters none of the early minerals. From the volatile-free \ce{K2O} of
lavas 1 and 2, compute the fraction $F$ of liquid left when lava 2 formed, if it
derives from lava 1 by \omterm{def:b2:magmatic-differentiation:crystallisation}{fractional crystallisation}.
\item Do the same with rubidium and with thorium.
\item Zirconium has a bulk $D \approx 0.1$. Use the Rayleigh law to estimate $F$ from
it.
\item Strontium falls from lava 2 to lava 3. Compute $F$ from potash between lavas 2 and
3, and use it to estimate the bulk $D$ of Sr over that interval. Which mineral is
responsible?
\item From lava 1 to lava 3, what fraction of the liquid crystallised?
\item If the \omterm{def:b2:magma-emplacement:magma}{magma} had assimilated granitic crust rich in potash, would the potash
estimate be too high or too low?
\item Average the four estimates of Part IV for lavas 1 to 2 and state the mass of
crystals left behind per \qty{100}{kg} of basalt to make the andesite, the key number
of this problem.
\end{enumerate}
\end{problem}
% ledger: bgs:tas-use, bgs:tas, bgs:k2o-sio2, geoplotr:afm, geoplotr:tizry, ms95:pm-trace
